\documentclass[11pt]{article}
\usepackage[margin=1in]{geometry}
\usepackage{amsmath,amssymb}
\usepackage{booktabs}
\usepackage{hyperref}

\title{CMB Constraint Mapping for Dark-Sector Necessity in Tau Core}
\author{Tau Core Research}
\date{Private scaffold, June 2026}

\begin{document}
\maketitle

\begin{abstract}
This private scaffold defines a CMB-facing constraint map for deciding what a
Tau Core cosmology must reproduce before any dark-sector necessity claim can be
made. The paper does not fit CMB spectra, does not implement a Boltzmann solver,
and does not claim that Tau Core explains the CMB. It instead maps CMB acoustic,
lensing, damping, and growth constraints to minimum Tau Core deliverables. The
main provisional conclusion is methodological: a completely no-dark Tau Core
cosmology is not a defensible global claim unless it can reproduce the CMB roles
normally carried by non-baryonic gravitational sources in standard cosmology.
\end{abstract}

\section{Scope}

Paper 13 is separated from the galaxy-rotation scoring work. Galaxy rotation can
test morphology/readout behavior, but the cosmic microwave background (CMB)
tests early-universe acoustic physics, gravitational potentials, lensing, and
growth seeds. Therefore the CMB is a stronger arena for asking whether Tau Core
can remain no-dark or whether an effective base-field response layer is needed.

This scaffold distinguishes three possible later outcomes:
\[
\begin{array}{ll}
\mathrm{TC0}: & \hbox{no effective dark-sector response required},\\
\mathrm{TC1}: & \hbox{Tau base-field response required},\\
\mathrm{TC2}: & \hbox{standard CDM/Lambda-like benchmark still required}.
\end{array}
\]
None of these outcomes is claimed here.

\section{External Constraints}

The primary reference point is the Planck 2018 cosmological parameter analysis,
which reports strong consistency of the base six-parameter flat
$\Lambda$CDM model with CMB anisotropy and lensing data
\cite{planck2018parameters}. The relevant pressure on Tau Core is not the name
``dark matter'' but the role: early gravitational potentials, acoustic peak
structure, matter-density inference, lensing, and growth consistency.

ACT DR6 CMB lensing provides an independent lensing-based check on the matter
distribution and growth history \cite{actdr6lensing}, while ACT DR6 extended
model analyses provide a current reference for how CMB data constrain departures
from base $\Lambda$CDM \cite{actdr6extended}.

\section{Tau Core Requirement Map}

The constraint map is intentionally conservative. A no-dark Tau Core cosmology
must not merely fit a late-time galaxy rotation curve. It must also provide:

\begin{enumerate}
\item an equation-level source model for the acoustic peak phase and ratios;
\item an effective-density mapping for the role conventionally captured by
  $\Omega_c h^2$;
\item a lensing-transfer response consistent with CMB lensing;
\item a radiation-content and damping-tail audit;
\item a growth-transfer kernel connecting early perturbations to later
  structure.
\end{enumerate}

If these roles cannot be delivered without an effective non-baryonic source,
then Tau Core must be formulated with a base-field response layer before making
cosmological claims. Such a layer should not be overpromoted to a detected
physical field; it is first a required effective role in the theory.

\section{Claim Boundary}

This paper does not claim that Tau Core explains the CMB, does not claim that
dark matter does not exist, does not detect a Tau dark-sector field, and does
not provide a CMB likelihood. It provides only a reproducible constraint map and
a next-artifact target:
\[
\texttt{tau\_cmb\_effective\_source\_equation\_v01}.
\]

\section{Reproducibility}

The repository includes compact CSV manifests and a script:
\[
\texttt{python scripts/reproduce.py}.
\]
The script regenerates derived tables, validates claim boundaries, builds the
arXiv source package, and runs public tests.

\bibliographystyle{unsrt}
\bibliography{refs}
\end{document}
